\section{Aljabar Abstrak Dasar untuk Kriptografi}

Kriptografi modern, terutama algoritma seperti AES dan ECC, tidak hanya menggunakan aritmetika bilangan bulat biasa, tetapi menggunakan struktur aljabar yang lebih luas. Penggunaan aljabar abstrak memungkinkan kriptografer untuk mendefinisikan operasi yang memiliki sifat matematis terjamin, seperti keberadaan invers untuk setiap elemen, yang sangat penting untuk proses dekripsi.

\subsection{Struktur Dasar: Grup, Ring, dan Field}

Dalam aljabar abstrak, kita mengelompokkan himpunan yang dilengkapi dengan operasi tertentu ke dalam beberapa struktur:

\begin{enumerate}
    \item \textbf{Grup}: Sebuah himpunan $G$ dengan operasi biner $*$ disebut grup jika memenuhi empat syarat:
    \begin{itemize}
        \item \textbf{Tertutup}: Untuk setiap $a, b \in G$, maka $a * b \in G$.
        \item \textbf{Asosiatif}: $(a * b) * c = a * (b * c)$.
        \item \textbf{Elemen Identitas}: Terdapat $e \in G$ sehingga $a * e = e * a = a$.
        \item \textbf{Elemen Invers}: Untuk setiap $a \in G$, terdapat $a^{-1} \in G$ sehingga $a * a^{-1} = e$.
    \end{itemize}
    Contoh: Himpunan bilangan bulat $\mathbb{Z}$ dengan operasi penjumlahan adalah sebuah grup.

    \item \textbf{Ring}: Sebuah himpunan yang memiliki dua operasi (biasanya penjumlahan dan perkalian). Penjumlahan harus membentuk grup abelian, sedangkan perkalian harus asosiatif dan distributif terhadap penjumlahan.

    \item \textbf{Field (Lapangan)}: Ring yang paling kuat. Dalam sebuah Field, setiap elemen non-nol harus memiliki invers multiplikatif. Artinya, di dalam Field, kita dapat melakukan semua operasi aritmetika dasar: penjumlahan, pengurangan, perkalian, dan pembagian (kecuali pembagian dengan nol).
\end{enumerate}

\subsection{Galois Field $\mathrm{GF}(2^n)$}

Dalam kriptografi digital, kita bekerja dengan bit. Oleh karena itu, kita menggunakan \textbf{Finite Field} atau \textbf{Galois Field}, yaitu field yang hanya memiliki jumlah elemen terbatas. Field yang paling dominan dalam kriptografi adalah $\mathrm{GF}(2^n)$.

Jenis yang paling umum digunakan dalam AES adalah $\mathrm{GF}(2^8)$, yaitu field dengan $2^8 = 256$ elemen. Setiap elemen dalam $\mathrm{GF}(2^8)$ direpresentasikan sebagai satu byte (8 bit), yang secara matematis dipandang sebagai polinomial derajat maksimal 7 dengan koefisien $\{0, 1\}$.

\textbf{Operasi dalam $\mathrm{GF}(2^8)$}:
\begin{itemize}
    \item \textbf{Penjumlahan}: Dilakukan menggunakan operasi XOR bitwise. Hal ini sangat efisien di tingkat perangkat keras komputer.
    \item \textbf{Perkalian}: Dilakukan dengan mengalikan dua polinomial, kemudian mengambil sisa baginya terhadap sebuah polinomial irreducible (polinomial yang tidak dapat difaktorkan), mirip dengan konsep modulo pada bilangan bulat. Untuk AES, polinomial yang digunakan adalah $m(x) = x^8 + x^4 + x^3 + x + 1$.
\end{itemize}

\subsection{Mengapa Menggunakan Aljabar Abstrak?}

Penggunaan $\mathrm{GF}(2^8)$ pada AES memastikan beberapa hal penting:
\begin{enumerate}
    \item \textbf{Keberadaan Invers}: Setiap byte non-nol memiliki invers unik dalam $\mathrm{GF}(2^8)$. Hal ini memungkinkan pembuatan \textit{S-Box} (Substitution Box) yang bersifat satu-ke-satu, sehingga proses dekripsi dapat dilakukan secara konsisten.
    \item \textbf{Efisiensi Komputasi}: Operasi XOR dan perkalian polinomial dapat diimplementasikan dengan sangat cepat menggunakan instruksi CPU modern.
    \item \textbf{Difusi Data}: Operasi dalam Galois Field memastikan bahwa perubahan satu bit pada plainteks akan menyebar ke banyak bit pada cipherteks (efek avalanche).
\end{enumerate}
